% GENERATED FILE. Edit canonical lessons, then run npm run generate:books.
% canonical-source-hash: fnv1a64:ca04cca9f23e142c
% canonical-lessons: FR-C306-etouffer, FR-C306-intoxication, FR-C306-piquer, FR-C306-mordre, FR-C306-morsure

\chapter{Petits accidents}
\label{ch:fr-petits-accidents}
\hlchaptermodality{306}

\begin{chapteropening}
\textbf{By the end of this chapter:} \emph{I can describe small accidents: choking, food poisoning, a sting, a bite and biting.}

Complete the last lesson of chapter 306: I can describe small accidents: choking, food poisoning, a sting, a bite and biting.
\end{chapteropening}

\section[étouffer]{étouffer — to choke (as s'étouffer); to suffocate, to stifle}

\label{lesson:FR-C306-etouffer}



\begin{quote}
Before the new one: say the French for to survive, then the French for the catastrophe.
\end{quote}

\subsection*{You'll want to know: étouffer}
\textbf{étouffer} — \textquotedblleft{}to choke (as s'étouffer); to suffocate, to stifle\textquotedblright{}.

\begin{quote}
Il s'est étouffé avec un bonbon.
\end{quote}

\begin{quote}
\emph{He choked on a sweet.}
\end{quote}

For choking on food, the reflexive \emph{s'étouffer}. Without \emph{se}, to smother or to stifle: \emph{on étouffe ici !} — it is stifling in here.

\begin{cousinweb}
\emph{étouffer} comes from Old French \emph{estouffer} (EtymDB); the origin of the root is not settled.
\end{cousinweb}

\begin{grammarlens}[title={What you've built}]
One more word for telling, explaining and describing at B1.
\end{grammarlens}

\subsection*{Your turn}
\begin{itemize}
\raggedright
  \item \emph{Say it:} \emph{étouffer}
  \item \emph{Say it:} \emph{étouffer}, once more
  \item \emph{Say it:} a sentence of your own with \emph{étouffer}
  \item \emph{Recall:} say the French for to survive, then the French for the catastrophe, then say \emph{étouffer} again
\end{itemize}

\begin{tcolorbox}[breakable,colback=teal!4,colframe=teal!35!black,title={Before you move on}]
What does étouffer mean? (\textquotedblleft{}To choke (as s'étouffer); to suffocate, to stifle\textquotedblright{}.) Say it once more.
\end{tcolorbox}
\section[l'intoxication]{l'intoxication — the poisoning}

\label{lesson:FR-C306-intoxication}



\begin{quote}
Before the new one: say the French for the catastrophe, then the French for to choke.
\end{quote}

\subsection*{You'll want to know: l'intoxication}
\textbf{l'intoxication} — \textquotedblleft{}the poisoning\textquotedblright{}.

\begin{quote}
Toute la famille a eu une intoxication alimentaire.
\end{quote}

\begin{quote}
\emph{The whole family had food poisoning.}
\end{quote}

\emph{Une intoxication alimentaire} — food poisoning. It does not mean being drunk.

\begin{cousinweb}
\emph{intoxication} is built on \emph{intoxiquer}, from Medieval Latin \emph{intoxicāre} (EtymDB), built on \emph{toxicum} \textquotedblleft{}poison\textquotedblright{}, a borrowing from Greek \emph{toxikón} \textquotedblleft{}poison for arrows\textquotedblright{}.
\end{cousinweb}

\begin{grammarlens}[title={What you've built}]
One more word for telling, explaining and describing at B1.
\end{grammarlens}

\subsection*{Your turn}
\begin{itemize}
\raggedright
  \item \emph{Say it:} \emph{l'intoxication}
  \item \emph{Say it:} \emph{l'intoxication}, once more
  \item \emph{Say it:} a sentence of your own with \emph{l'intoxication}
  \item \emph{Recall:} say the French for the catastrophe, then the French for to choke, then say \emph{l'intoxication} again
\end{itemize}

\begin{tcolorbox}[breakable,colback=teal!4,colframe=teal!35!black,title={Before you move on}]
What does l'intoxication mean? (\textquotedblleft{}The poisoning\textquotedblright{}.) Say it once more.
\end{tcolorbox}
\section[piquer]{piquer — to sting, to bite (of insects); to prick}

\label{lesson:FR-C306-piquer}



\begin{quote}
Before the new one: say the French for to choke, then the French for the poisoning.
\end{quote}

\subsection*{You'll want to know: piquer}
\textbf{piquer} — \textquotedblleft{}to sting, to bite (of insects); to prick\textquotedblright{}.

\begin{quote}
Un insecte l'a piqué au bras.
\end{quote}

\begin{quote}
\emph{An insect stung him on the arm.}
\end{quote}

Informal: \emph{piquer} also means to pinch, to steal.

\begin{cousinweb}
\emph{piquer} is inherited from Vulgar Latin \emph{piccāre} \textquotedblleft{}to prick, to sting\textquotedblright{} (EtymDB).
\end{cousinweb}

\begin{grammarlens}[title={What you've built}]
One more word for telling, explaining and describing at B1.
\end{grammarlens}

\subsection*{Your turn}
\begin{itemize}
\raggedright
  \item \emph{Say it:} \emph{piquer}
  \item \emph{Say it:} \emph{piquer}, once more
  \item \emph{Say it:} a sentence of your own with \emph{piquer}
  \item \emph{Recall:} say the French for to choke, then the French for the poisoning, then say \emph{piquer} again
\end{itemize}

\begin{tcolorbox}[breakable,colback=teal!4,colframe=teal!35!black,title={Before you move on}]
What does piquer mean? (\textquotedblleft{}To sting, to bite (of insects); to prick\textquotedblright{}.) Say it once more.
\end{tcolorbox}
\section[mordre]{mordre — to bite}

\label{lesson:FR-C306-mordre}



\begin{quote}
Before the new one: say the French for the poisoning, then the French for to sting.
\end{quote}

\subsection*{You'll want to know: mordre}
\textbf{mordre} — \textquotedblleft{}to bite\textquotedblright{}.

\begin{quote}
Le chien du voisin a mordu le facteur.
\end{quote}

\begin{quote}
\emph{The neighbour's dog bit the postman.}
\end{quote}

\emph{Se faire mordre} — to get bitten.

\begin{cousinweb}
\emph{mordre} is inherited from Latin \emph{mordēre} \textquotedblleft{}to bite\textquotedblright{} (EtymDB).
\end{cousinweb}

\begin{grammarlens}[title={What you've built}]
One more word for telling, explaining and describing at B1.
\end{grammarlens}

\subsection*{Your turn}
\begin{itemize}
\raggedright
  \item \emph{Say it:} \emph{mordre}
  \item \emph{Say it:} \emph{mordre}, once more
  \item \emph{Say it:} a sentence of your own with \emph{mordre}
  \item \emph{Recall:} say the French for the poisoning, then the French for to sting, then say \emph{mordre} again
\end{itemize}

\begin{tcolorbox}[breakable,colback=teal!4,colframe=teal!35!black,title={Before you move on}]
What does mordre mean? (\textquotedblleft{}To bite\textquotedblright{}.) Say it once more.
\end{tcolorbox}
\section[la morsure]{la morsure — the bite (a wound)}

\label{lesson:FR-C306-morsure}



\begin{quote}
Before the new one: say the French for to sting, then the French for to bite.
\end{quote}

\subsection*{You'll want to know: la morsure}
\textbf{la morsure} — \textquotedblleft{}the bite (a wound)\textquotedblright{}.

\begin{quote}
La morsure du chien n'était pas grave.
\end{quote}

\begin{quote}
\emph{The dog bite was not serious.}
\end{quote}

For insects, French uses the verb \emph{piquer}.

\begin{cousinweb}
\emph{morsure} is built on \emph{mors}, the old past participle of \emph{mordre} \textquotedblleft{}to bite\textquotedblright{}, from Latin \emph{morsus} \textquotedblleft{}a bite\textquotedblright{} (Lewis \& Short).
\end{cousinweb}

\begin{grammarlens}[title={What you've built}]
One more word for telling, explaining and describing at B1.
\end{grammarlens}

\subsection*{Your turn}
\begin{itemize}
\raggedright
  \item \emph{Say it:} \emph{la morsure}
  \item \emph{Say it:} \emph{la morsure}, once more
  \item \emph{Say it:} a sentence of your own with \emph{la morsure}
  \item \emph{Recall:} say the French for to sting, then the French for to bite, then say \emph{la morsure} again
\end{itemize}

\begin{tcolorbox}[breakable,colback=teal!4,colframe=teal!35!black,title={Before you move on}]
What does la morsure mean? (\textquotedblleft{}The bite (a wound)\textquotedblright{}.) Say it once more.
\end{tcolorbox}
